\section{History-gap and source calculations}
\label{app:source-details}

\subsection{Rational gain/loss split}
\label{app:rational-split}

On a clock bit \(c\) and work bit \(w\), define
\[
\begin{aligned}
u_0&=(\ket{00}+\ket{01}-\ket{10})/\sqrt3,&
u_1&=(\ket{00}-\ket{01}-\ket{11})/\sqrt3,\\
d_0&=(\ket{00}-\ket{10}-\ket{11})/\sqrt3,&
d_1&=(\ket{01}-\ket{10}+\ket{11})/\sqrt3.
\end{aligned}
\]
The vectors in each pair are orthonormal.  If consecutive history amplitudes
are \(f_t,f_{t+1}\), orthogonality to \(u_0,u_1\) is equivalent to
\(f_{t+1}=\sqrt2Hf_t\); orthogonality to \(d_0,d_1\) is equivalent to
\(f_{t+1}=Hf_t/\sqrt2\).  Applying the first relation to one work bit and the
second to the other gives \(H\otimes H\).

\subsection{Weighted path estimate}

With the unary encoding \(\ket t=\ket{1^t0^{T-t}}\), the legal clock span and
its orthogonal complement are invariant under the neighboring-clock-guarded
propagation, input, and output projectors.  The diagonal clock term
\[
H_{\rm clock}=\sum_{j=1}^{T-1}\ket{01}\!\bra{01}_{j,j+1}
\]
has energy at least one on every illegal computational-basis clock string.
By Hermiticity, the whole unary-clock Hilbert space is the orthogonal sum of
the invariant legal and illegal sectors.  It therefore
suffices to prove the smaller inverse-polynomial bound on the legal sector.
For a legal-clock state, write \(f_t=s_tV_t\xi_t\), with the notation of~\cref{eq:history-isometry}. Its norm and initial-history energy are
\[
N(\xi)=\sum_tw_t\|\xi_t\|^2,
\qquad
E(\xi)=\sum_{t=1}^{L-1}c_t\|\xi_t-\xi_{t-1}\|^2
+\langle\xi_0,A_{\rm in}\xi_0\rangle,
\]
where \(w_t=s_t^2\in[1,2]\), \(c_t\in\{1/2,2/3\}\), and
\(A_{\rm in}\) is the clean-input penalty at time zero, a sum of commuting
computational-basis projectors, one per clean qubit.  It is not in general a
projector: a basis string violating \(k\) clean bits has energy \(k\).  All
that is used below is
\[
\ker A_{\rm in}=C,\qquad A_{\rm in}\succeq P_{C^\perp}.
\]
Split every \(\xi_t\) into the clean subspace \(C\) and its orthogonal
complement, and split the energy accordingly as \(E(\xi)=E_C+E_B\) with
\[
E_C=\sum_{t=1}^{L-1}c_t\|\xi_t^C-\xi_{t-1}^C\|^2,\qquad
E_B=\sum_{t=1}^{L-1}c_t\|\xi_t^{C^\perp}-\xi_{t-1}^{C^\perp}\|^2
+\langle\xi_0,A_{\rm in}\xi_0\rangle ;
\]
the two difference terms are orthogonal because \(C\) is invariant under
the splitting.

For a state orthogonal to all valid histories, the clean component has weighted mean zero:
\[
\sum_tw_t\xi_t^C=0.
\]
The weighted variance identity and path Cauchy--Schwarz inequality give
\[
\sum_tw_t\|\xi_t^C\|^2
=\frac1{2\sum_tw_t}
\sum_{i,j}w_iw_j\|\xi_i^C-\xi_j^C\|^2
\le2L^2E_C.
\]
For the dirty component, set
\(a=\|\xi_0^{C^\perp}\|^2\le\langle\xi_0,A_{\rm in}\xi_0\rangle\le E_B\) and
\(
D_B=\sum_t\|\xi_t^{C^\perp}-\xi_{t-1}^{C^\perp}\|^2\le2E_B
\), using \(c_t\ge1/2\).
The bound
\[
\|\xi_t^{C^\perp}\|^2\le2a+2(L-1)D_B
\]
and \(\sum_tw_t\le2L\) imply
\[
\sum_tw_t\|\xi_t^{C^\perp}\|^2
\le2L\bigl(2E_B+4(L-1)E_B\bigr)\le8L^2E_B.
\]
Adding both sectors proves \(N(\xi)\le8L^2E(\xi)\) on the orthogonal complement of the complete initial kernel, establishing the first bound in~\cref{eq:history-gaps}.

\subsection{Adding the output projector}

We record the two-positive-operator estimate used for the second bound in~\cref{eq:history-gaps}. Let \(P\) project onto \(\ker A\), suppose
\(A\succeq g(I-P)\), and let \(R\) be an orthogonal projector.  Put
\(\mathcal S=\ker A\cap\ker R\), and suppose \(PRP\succeq\alpha P\) on
\(\ker A\ominus\mathcal S\).  On every nontrivial canonical block for the
pair \(P,R\), for some \(\rho\in[\alpha,1]\),
\[
 P=\begin{pmatrix}1&0\\0&0\end{pmatrix},\qquad
 R=\begin{pmatrix}
 \rho&\sqrt{\rho(1-\rho)}\\
 \sqrt{\rho(1-\rho)}&1-\rho
 \end{pmatrix}.
\]
Thus \(g(I-P)+R\) has trace \(g+1\), determinant \(g\rho\), and
smallest eigenvalue
\[
 \frac{g+1-\sqrt{(g+1)^2-4g\rho}}2.
\]
This quantity increases with \(\rho\); the one-dimensional canonical blocks
outside \(\mathcal S\) obey the same lower bound.  Since
\(A+R\succeq g(I-P)+R\) and both operators have kernel \(\mathcal S\),
\begin{equation}
\label{eq:two-positive}
\gamma_+(A+R)
\ge
\frac{g+1-\sqrt{(g+1)^2-4g\alpha}}2
\ge\frac{g\alpha}{g+1}.
\end{equation}
For the history Hamiltonian, \(g=1/(8L^2)\) and \(\alpha=(1-r)/Z\ge1/(3L)\) when \(r\le1/3\). Substitution into~\cref{eq:two-positive} gives
\[
\gamma_+(H_{\rm out})
\ge\frac1{3L(8L^2+1)}.
\]
The kernel identity follows directly from positivity:
\[
\ker(H_{\rm in}+H_{\rm rej})
=\ker H_{\rm in}\cap\ker H_{\rm rej}
=\cV\ker(I-M).
\]

\subsection{Exact eight-label compression}

Let \(U_x\) denote the fixed-input verifier and \(q\) its output bit. The label \(z\) is never changed. After computing and uncomputing the constant-size acceptance predicate, compression through all clean work registers has no matrix element between distinct labels. On \(z=000\), acceptance is the identity. On five labels it is the scalar
\(
p_x=\|P_{\rm acc}U_x\ket{0\cdots0}\|^2
\),
and on two labels it is zero. This proves~\cref{eq:eight-label} even for a coherent vector across label values.

The normalized trace is
\[
\frac{\operatorname{Tr}M_x}{8}
=\frac{1+5p_x}{8}.
\]
If \(p_x=1\), this equals \(3/4\) and coincides with the perfect-space fraction. If \(p_x\le1/3\), it is at most \(1/3\), the perfect-space fraction is exactly \(1/8\), and every off-perfect eigenvalue is at most \(1/3\). Thus the construction simultaneously supplies the trace promise, the perfect-space separation, and the off-perfect ceiling needed by the output-gap proof.

\subsection{Source-interface boundary}
\label{app:source-boundaries}

The exact source used by \cref{thm:weighted-hardness} begins with a
fixed-input \(\BQPone^{G_2}\) verifier.  The three mixed label bits are added
by the reduction and are the entire input space of
\cref{eq:eight-label}; the verifier's input and all predicate scratch remain
clean.  Consequently the initial history dimension is exactly eight.  No
trace-to-rank inference, unrestricted \(\SDQCone\) identification, or
approximate gate compilation enters the proof.

The finite gadget evidence needed to compile these circuit terms is given in
\cref{app:palette}.  The theorem uses those exact finite objects and does not
infer their topology from floating-point spectra.

For an arbitrary acceptance operator with
\(M|_{S^\perp}\preceq rI\), put \(f=\dim S/D\) and
\(p=\operatorname{Tr}(M)/D\).  Spectral decomposition gives
\[
f\le p\le r+(1-r)f.
\]
Thus trace promises \(p\ge a\) versus \(p\le b\) force a perfect-space
fraction gap only when
\(\max\{0,(a-r)/(1-r)\}>b\).  The eight-label source avoids this obstruction
by fixing the perfect fractions directly.

An exact circuit over \(\{X,\CX,\CCX,H\}\) admits the following real-gate
extension.  Add one maximally mixed, unmeasured spectator \(a\) and replace
each \(H_j\) by \(H_j\otimes H_a\).  If there are \(h\) Hadamards, then
\[
U'=U\otimes H_a^h,\qquad J'=J\otimes I_a,\qquad
P'_{\rm acc}=P_{\rm acc}\otimes I_a.
\]
For either parity of \(h\),
\[
(J')^*(U')^*P'_{\rm acc}U'J'=M\otimes I_a.
\]
Both the input dimension and perfect eigenspace dimension double, while the
fraction and off-perfect ceiling are unchanged.  This calculation does not
cover complex phases.
